\section{Reference families by a dominated lifted coupling}\label{sec:reference}
We first distinguish two uses of DH. Its Section 3 averaging construction,
Propositions 5.1--5.2, Theorem 5.3 and Corollary 4.7 apply under A2$'$:
translation ergodicity, square symmetries, a $2+\eps$ moment, bounded norm
shape, and unique passage times. Lemmas~\ref{lem:capfacts}--\ref{lem:metric}
supply these for $s$. DH Section 6, including Lemma 6.2 and the following
triple extraction, has a \emph{standing upward-finite-energy assumption}.
We reuse the proofs of those geometric extraction steps, which use only
supporting-arc geometry, stationarity and the ergodic theorem; we do not
invoke their statements under an unavailable hypothesis. The later
finite-energy surgery is replaced below.

Use DH's exact rounded-line targets
\[
 \widehat L_\alpha=\{z\in\Z^2:(z+[-1/2,1/2)^2)
                  \cap\{x:\rho(x)=\alpha\}\ne\varnothing\}.
\]
No identification of line-target and half-plane-target graph laws is needed.
Coercivity and uniqueness give finite geodesics to each such target. The DH
averages, lifted to include $X$, $s$, graph and cocycle coordinates, have a
stationary subsequential limit $\nu$. Its paths are geodesics, have outdegree
one, merge on intersection, and lie in the supporting arc of the limiting
cocycle slope. The geometric proof of DH Lemma 6.2 and its subsequent
extraction produce, on noncoalescence, a positive-probability triple of
disjoint rays eventually in every translate of one of eight half-planes.
Ergodicity of the lifted law is not needed: its relevant ergodic average
may be random with positive expectation.

\begin{lemma}[D: domination and selected-ray stability]\label{lem:lift}
Fix finite $F\subset\Ecal$ and $a$ with $r=G((a,\infty))>0$.
Resample $t_F$ independently from $G(\cdot\mid t>a)$, keeping everything
else, including $V$, fixed; denote the resulting field by $X'$. There is
a joint limiting coupling with old marginal $\nu$ and new marginal
$\nu'\le r^{-|F|}\nu$. Its underlying pair marginal is exactly this
original-coordinate coupling. On $\max_Ft<a$, any entirely uncapped old
limiting ray avoiding $F$ is the selected new limiting ray from the same root.
\end{lemma}
\begin{proof}
The new original marginal has density
$r^{-|F|}\prod_{e\in F}\ind_{\{t'_e>a\}}\le r^{-|F|}$ relative to the
old iid law. Recompute every finite-line graph and cocycle with the same
$\alpha$ in the two fields. The density bound survives these pushforwards
and averaging over $\alpha$. Joint tightness follows from the fixed
underlying pair law, compact graph coordinates, and cocycle increments
bounded by cap distances; the new distance laws are dominated by a fixed
multiple of the old ones. Take a joint subsequential limit. Domination
passes first for nonnegative bounded continuous functions, then for all
Borel sets by regularity.

Cap coordinates agree with the cap of the retained originals in the limit.
Indeed a shortest circuit and all potentially shorter circuits depend on
finitely many indicators. None of their originals equals $M$ almost surely,
so $L_e$ is locally constant and $s_e$ continuous at almost every original
configuration. Absolute continuity supplies the same fact for the new marginal.

For selected-ray stability fix a root $v$, a finite desired prefix and
$\eps>0$. By finite influence choose $R$ containing $F,v$ so the cap-change
set lies in $Q_R$ with probability at least $1-\eps$. Choose $S\ge R$ so
all old geodesics between vertices of $Q_R$ stay in $Q_S$ with probability
at least $1-\eps$. There are finitely many endpoint pairs; coercivity gives
such an $S$. Both failure bounds are uniform over the approximants because
their underlying marginal is fixed. On this event a path from $v$ cannot
return to $Q_R$ after leaving $Q_S$.

Realize the joint convergence by Skorokhod. Every edge of the asserted old
ray satisfies the strict inequality $t_e<K_e$ almost surely; equality has
probability zero by the conditional nonatomic variables. Its finite prefix
through the first exit from $Q_S$ agrees in late old graph approximants,
and the strict margins and avoidance of $F$ persist. On $\max_Ft<a$ the
change is upward. This prefix is unchanged; on the two localization events
the rest of the finite target path is outside $Q_R$ and unchanged as well.
The whole finite path remains the unique new minimizer to the same target.

If the limiting desired prefixes differed, they would eventually differ
in the discrete approximants. On the asserted limiting event this forces
a localization failure in every sufficiently late approximant. Fatou bounds
the bad limiting probability by $2\eps$. Let $\eps\downarrow0$ and intersect
over roots and prefix lengths. This does not assume almost-sure eventual
localization in the Skorokhod realization.
\end{proof}

\begin{remark}\label{rem:kernels}
The same proof applies to a kernel changing a fixed deterministic finite
set of original coordinates whose
output marginal has a bounded density relative to the underlying product
law, on an event where that kernel increases the originals. The joint
underlying law is fixed across the approximants, the density bound gives
the same domination, and strict uncapping plus the two tight localization
events give the identical stability conclusion. If its upward event is
only measurable, include that event's indicator in the joint lift. Its
discrete convergence makes the event persist in the approximants, where
monotonicity holds; the entire underlying pair-and-indicator marginal is
fixed at every approximation and hence in the limit. In particular,
same-status resampling has output density one, while each finite tilt has
a bounded product density. This extends the explicit reference and label
couplings below to those kernels; it does not assert cap conditional finite
energy or stability under arbitrary uncontrolled modifications. The later
edges and regions are deterministic once their site and radius are fixed;
there are only countably many such choices. A random support is not being
covered by an unstated uniform kernel bound.
\end{remark}

The line convention is compatible with the circuit shortcut: if a rounded
target vertex lies strictly inside a lattice circuit, its half-open unit
cell is inside the polygon. The line enters the polygon and exits at a
boundary edge or vertex, where rounding gives a boundary target vertex.
A cheap boundary arc therefore reaches a valid target. This also justifies
the uncapping argument for the precise DH approximants.

\subsection{W: the cubic wall and its diagonal version}\label{sec:wall}
Suppose the limiting graph is not coalescing. Cut the extracted triple at
last contacts with a sufficiently far common level, after all capped edges.
Countably many choices of roots, level and orientation, followed by
stationarity, give a deterministic positive-probability event of three
wholly uncapped rays rooted on that level and strictly to its positive side
afterward. Selector consistency follows from outdegree one and merging.

For an axis level take roots $(0,j_1),(0,j_2),(0,j_3)$ with
$0<j_1<j_2<j_3<m$, and let $F_m$ consist of the rungs
$(-1,j)(0,j)$, $1\le j\le m-1$. Raise their originals above
$H_m=4Mm^3$. A protecting open circuit enclosing any rung has crossings
of $x=-1/2$ on both sides of it. The intermediate crossings are all closed,
so the vertical span is at least $m$. Hence $L'_e\ge2m$ and
$K'_e\ge8Mm^3>H_m$; the new rung cap is larger than $H_m$.

With probability tending to one, an open circuit strictly surrounding this
wall has at most $C_0m^2$ edges. It avoids the raised set and costs at most
$C_0Mm^2<H_m$ for large $m$. Choose such a deterministic $m$ with positive
joint probability of the triple and cheap circuit, then choose finite
$a\ge H_m$ so that $\max_{F_m}t<a$ still leaves positive probability.
Resample above $a$ as in Lemma~\ref{lem:lift}. The three selected rays persist.

The outer two rays and the level segment joining their roots separate the
plane. A vertex outside the cheap circuit and outside the right-hand corridor
cannot be in the middle class: a path to a sufficiently late middle-ray
vertex would meet an outer ray, merging with a different class, or cross a
rung. In the latter case its segment between two contacts with the cheap
circuit can be bypassed at smaller cost. Thus the middle class lies in the
corridor plus a finite set. It is bounded on the left and meets every fixed
vertical column finitely, since the bounding rays eventually lie beyond
every such column. It has a least vertex in lexicographic $(x,y)$ order.

For a diagonal level use $h=x+y$, $u=x-y$, with roots $(k,-k)$ on $h=0$.
Raise both incoming edges from $h=-1$ into each $(k,-k)$, $|k|\le m$.
Their crossings of $h=-1/2$ have $u=2k-1/2$ and $u=2k+1/2$: consecutive
half-integers from $-2m-1/2$ to $2m+1/2$. A protecting open circuit must
cross on both sides of this blocked interval; its $u$-span exceeds $4m$,
so $L'_e\ge4m$ and $K'_e\ge64Mm^3>H_m$. The same $H_m=4Mm^3$ pays the
$O(Mm^2)$ outer circuit. Entry into the corridor through the interior of
the level segment uses one of the raised incoming edges; entry at an outer
root meets an outer ray. Lexicographic $(h,u)$ order gives the least vertex.
This parity calculation is explicit; a $45$-degree rotation is not being
treated as a square-lattice symmetry. Actual square rotations and reflections
cover the remaining six normals.

Under the stationary old graph law no class can have a least vertex in a
fixed translation-invariant order. Send one unit from each vertex to its
class's least vertex if one exists. Outgoing mass is at most one, whereas
that least vertex receives infinitely many units from its forward ray.
The mass-transport identity and countability rule this out. Domination
$\nu'\le r^{-|F|}\nu$ transfers the null event to the new law, contradicting
the construction. The limiting graph therefore coalesces.
This is the replacement of DH Section 6 surgery and, later, of AH 9.10
Case 1; no cap finite energy has been invoked. The ceiling condition of
DH Lemma 6.4 is vacuous here because the cap has unbounded essential support.

\subsection{Deterministic reference families and the explicit coupling}
For \emph{each tangent functional} $\rho$, choose a point of $\partial\{\mu\le1\}$
at which $\{\rho=1\}$ is the unique supporting line. This is AH's meaning
of tangent; it is precisely the differentiability hypothesis of DH
Corollary 4.7. That result makes the limiting cocycle slope deterministic
$\rho$. Once the graph coalesces its calibrated cocycle is its actual
Busemann function, by evaluating at a common future vertex.

We reuse the proofs of AH 5.3--5.6, not its finite-energy-dependent existence
statement: conditional quantiles of the ordered ray law give, for almost
every fixed quantile, a measurable translation-covariant coalescing family
of slope $\rho$. The already proved graph coalescence supplies the required
premise. The proof of AH 4.4 identifies any two such families. Its boundary
expectation step (AH 4.6) applies to the cap because coordinate, or diagonal
zigzag, boundary routes have finite mean and boundary translations are
ergodic. S's boundary-comparison appendix is unnecessary.

Fix one good quantile. Uniqueness for every other good fixed quantile,
followed by Fubini, says that the conditional ray law given $s$ is a point
mass. Countability of roots identifies the entire graph as a measurable
family $G_\rho(s)$. This statement is conditional on $s$, not merely on
the richer $X$. Domination transfers it to the modified graph.
Apply Lemma~\ref{lem:lift} at late vertices beyond capped edges and $F$.
The old and new $G_\rho$ rays share tails. Since both graphs are now
functions of their cap fields and the projected pair marginal is the
prescribed law of $(X,X')$, this stability holds under that \emph{explicit
original coupling}, not only along a chosen subsequence. Four square
rotations provide the reference families needed for the AH order construction.
